\chapter{Kommentierte Quellen}

Achtundzwanzig Werke, eingeordnet: Zu jedem steht, wofür man es aufschlägt und was es leistet, das andere nicht leisten.

\lit{1}{Anderson, R.: \textit{Security Engineering}. 3.\ Aufl., Wiley 2020.}{\textbf{Das Buch über alles, was in diesem Band nicht steht.} Anderson behandelt Sicherheit als Systemfrage: Menschen, Anreize, Organisationen, Seitenkanäle, Fehlbedienung. Über tausend Seiten mit unzähligen Fallbeispielen, frei verfügbar. Wer nach Kapitel~15 wissen will, warum Systeme tatsächlich fallen, liest hier weiter.}

\lit{2}{Arıkan, E.: \textit{Channel Polarization}. IEEE Transactions on Information Theory 55, 2009.}{Die Arbeit mit den Polar-Codes --- der erste explizite Code, der die Kapazität beweisbar erreicht, sechzig Jahre nach Shannon. Der Gedanke, gleiche Kanäle in ungleiche zu verwandeln, ist elementar und überraschend zugleich.}

\lit{3}{Beutelspacher, A.: \textit{Kryptologie}. 10.\ Aufl., Springer Spektrum 2015.}{\textbf{Auf Deutsch und für den Einstieg geeignet.} Beutelspacher erklärt die klassischen Verfahren und die Grundideen der modernen Kryptographie mit ungewöhnlich wenig Vorkenntnis. Wo dieser Band zu schnell war, ist das der sanftere Weg.}

\lit{4}{Boneh, D., Shoup, V.: \textit{A Graduate Course in Applied Cryptography}. Frei verfügbar, laufend aktualisiert.}{Die gründlichste zeitgemäße Darstellung mit Sicherheitsbeweisen. Behandelt genau die Fragen, die dieser Band nur streifen konnte: Was heißt es präzise, dass ein Verfahren sicher ist, und wie führt man den Beweis relativ zu einer Annahme?}

\lit{5}{Cover, T., Thomas, J.: \textit{Elements of Information Theory}. 2.\ Aufl., Wiley 2006.}{\textbf{Für die Teile I bis III die maßgebliche Quelle.} Enthält alle Beweise, die hier ausgelassen wurden --- insbesondere die vollständige Erreichbarkeit in Kapitel~7 und die Kraftsche Ungleichung für eindeutig lesbare Codes. Sauber geschrieben und mit reichlich Übungen.}

\lit{6}{Csiszár, I., Körner, J.: \textit{Information Theory: Coding Theorems for Discrete Memoryless Systems}. 2.\ Aufl., Cambridge 2011.}{Die theoretisch anspruchsvollste Darstellung. Entwickelt die Methode der Typen, mit der sich die Sätze aus Kapitel~5 mit exakten Fehlerexponenten statt nur asymptotisch beweisen lassen. Für Leser, die Kapitel~5 unbefriedigend ungenau fanden.}

\lit{7}{Diffie, W., Hellman, M.: \textit{New Directions in Cryptography}. IEEE Transactions on Information Theory 22, 1976.}{Die Arbeit, die alles veränderte. Enthält den Schlüsseltausch, den Begriff der Einwegfunktion und --- als Idee ohne Umsetzung --- die digitale Signatur. Lesenswert auch, weil die Verfasser die Tragweite selbst benennen.}

\lit{8}{Ferguson, N., Schneier, B., Kohno, T.: \textit{Cryptography Engineering}. Wiley 2010.}{\textbf{Das Buch über die Umsetzung.} Warum HMAC so aussieht, wie es aussieht; welche Betriebsart wann; wie man Zufall gewinnt; wie man Schlüssel verwaltet. Der Kapitel~15 zugrundeliegende Befund --- Verfahren fallen durch Umsetzung --- ist hier über dreihundert Seiten ausgeführt.}

\lit{9}{Gallager, R.: \textit{Information Theory and Reliable Communication}. Wiley 1968.}{Das klassische Lehrbuch des Fachs, von einem seiner wichtigsten Vertreter. Enthält die Zufallscodierungsschranke mit exakten Exponenten und ist bis heute die beste Quelle für die Frage, wie schnell die Fehlerrate tatsächlich fällt.}

\lit{10}{Goldreich, O.: \textit{Foundations of Cryptography}. Band I und II, Cambridge 2001/2004.}{Die begriffliche Grundlegung: Was ist eine Einwegfunktion, ein Pseudozufallsgenerator, ein Beweis ohne Wissenspreisgabe --- und wie hängt alles zusammen. Nicht für die Praxis, sondern für die Frage, worauf sie beruht.}

\lit{11}{Hankerson, D., Menezes, A., Vanstone, S.: \textit{Guide to Elliptic Curve Cryptography}. Springer 2004.}{Die Vertiefung zu Kapitel~14 --- einschließlich des Beweises, dass das Gruppengesetz assoziativ ist, und aller Umsetzungsfragen von der Punktdarstellung bis zur Abwehr von Seitenkanälen.}

\lit{12}{Huffman, D.: \textit{A Method for the Construction of Minimum-Redundancy Codes}. Proceedings of the IRE 40, 1952.}{Vier Seiten, die seit siebzig Jahren in Gebrauch sind. Der Optimalitätsbeweis ist kürzer als jede Nacherzählung und benutzt nichts als ein Vertauschungsargument.}

\lit{13}{Katz, J., Lindell, Y.: \textit{Introduction to Modern Cryptography}. 3.\ Aufl., Chapman \& Hall 2020.}{\textbf{Für die Teile IV und V das Standardlehrbuch.} Es unterscheidet sich von älteren Darstellungen dadurch, dass jedes Verfahren mit einer präzisen Sicherheitsdefinition und einem Beweis relativ zu einer benannten Annahme daherkommt --- genau die Sorgfalt, die Kapitel~13 als fehlend beklagte.}

\lit{14}{Kolmogorow, A.: \textit{Three approaches to the quantitative definition of information}. Problems of Information Transmission 1, 1965.}{Sieben Seiten, in denen der Begründer der Wahrscheinlichkeitstheorie einen Zufallsbegriff sucht, den seine eigene Theorie nicht liefern kann. Methodisch bemerkenswert und gut lesbar.}

\lit{15}{Li, M., Vitányi, P.: \textit{An Introduction to Kolmogorov Complexity and Its Applications}. 4.\ Aufl., Springer 2019.}{Das umfassende Werk zu Kapitel~7. Enthält die Inkomprimierbarkeitsmethode mit zahlreichen Anwendungen, Chaitins Unvollständigkeitsergebnisse und den Zusammenhang zur Lerntheorie. Über achthundert Seiten und die einzige vollständige Darstellung.}

\lit{16}{Lin, S., Costello, D.: \textit{Error Control Coding}. 2.\ Aufl., Pearson 2004.}{Das Standardwerk zu Kapitel~6. Behandelt lineare Codes, zyklische Codes, Reed-Solomon mit Decodierverfahren, Faltungscodes und LDPC --- mit allen Rechnungen, die hier nur angedeutet werden konnten.}

\lit{17}{MacKay, D.: \textit{Information Theory, Inference, and Learning Algorithms}. Cambridge 2003.}{\textbf{Das Buch, das die Brücke schlägt} --- zwischen Informationstheorie, Statistik und maschinellem Lernen, also zwischen den Teilen I bis III dieses Bandes und Band 17. Ungewöhnlich lebendig geschrieben, voller Beispiele und frei verfügbar. MacKay hat auch die LDPC-Codes wiederentdeckt.}

\lit{18}{Menezes, A., van Oorschot, P., Vanstone, S.: \textit{Handbook of Applied Cryptography}. CRC Press 1996.}{Das Nachschlagewerk. Kein Lehrbuch, sondern eine systematische Sammlung von Verfahren, Angriffen und Zahlenwerten, frei verfügbar. Wer wissen will, welche Varianten es gibt und was gegen sie bekannt ist, beginnt hier.}

\lit{19}{Nielsen, M., Chuang, I.: \textit{Quantum Computation and Quantum Information}. Cambridge 2010.}{Das Standardwerk zum Quantenrechnen. Enthält Shors und Grovers Verfahren vollständig sowie die Quantenfehlerkorrektur --- und einen eigenen Teil über Quanteninformationstheorie, in dem die Größen aus Teil~I ihre quantenmechanische Entsprechung finden.}

\lit{20}{Paar, C., Pelzl, J.: \textit{Kryptografie verständlich}. Springer Vieweg 2016.}{\textbf{Auf Deutsch, für Ingenieure geschrieben.} Behandelt AES, RSA, elliptische Kurven und Hashfunktionen mit durchgerechneten Beispielen und Blick auf die Umsetzung in Hardware. Die passende Ergänzung zu den Kapiteln~10 bis 14.}

\lit{21}{Richardson, T., Urbanke, R.: \textit{Modern Coding Theory}. Cambridge 2008.}{Die Theorie der Codes, die der Kapazität nahekommen --- LDPC, Nachrichtenaustausch bei der Decodierung, Dichteentwicklung. Erklärt, warum Codes mit schlechtem Minimalabstand besser sein können, und das ist der Bruch mit Kapitel~6.}

\lit{22}{Rissanen, J.: \textit{Modeling by shortest data description}. Automatica 14, 1978.}{Die Begründung des Prinzips der kürzesten Beschreibung. Bemerkenswert ist der Verzicht auf die Annahme, dass es ein wahres Modell gibt --- eine Position, die in der Statistik bis heute die Minderheit ist.}

\lit{23}{Rivest, R., Shamir, A., Adleman, L.: \textit{A Method for Obtaining Digital Signatures and Public-Key Cryptosystems}. Communications of the ACM 21, 1978.}{Sieben Seiten, die die digitale Welt tragen. Enthält das Verfahren, den Korrektheitsbeweis und die Preisaufgabe, deren Sicherheitsschätzung um vierzehn Größenordnungen danebenlag.}

\lit{24}{Shannon, C.: \textit{A Mathematical Theory of Communication}. Bell System Technical Journal 27, 1948.}{\textbf{Die Arbeit, die alles enthält.} Entropie, Quellencodierung, Kapazität, Codierungstheorem, stetige Kanäle --- gut fünfzig Seiten, frei verfügbar und erstaunlich gut lesbar. Wer nur eine Originalarbeit dieses Bandes liest, liest diese.}

\lit{25}{Shannon, C.: \textit{Communication Theory of Secrecy Systems}. Bell System Technical Journal 28, 1949.}{Die Begründung der mathematischen Kryptographie. Enthält Satz 9.3, die Unizitätsdistanz sowie die Begriffe Verwirrung und Verwischung, die den Entwurf von Blockchiffren bis heute leiten.}

\lit{26}{Shor, P.: \textit{Polynomial-Time Algorithms for Prime Factorization and Discrete Logarithms on a Quantum Computer}. SIAM Journal on Computing 26, 1997.}{Die Arbeit, die Teil~V ein Ablaufdatum gab. Der Kern ist die Quanten-Fouriertransformation zur Periodenbestimmung; der Rest ist klassische Zahlentheorie. Bemerkenswert zugänglich für ein Ergebnis dieser Tragweite.}

\lit{27}{Werner, M.: \textit{Information und Codierung}. 3.\ Aufl., Springer Vieweg 2008.}{\textbf{Auf Deutsch, knapp und anwendungsnah.} Deckt die Teile I und II ab, mit Schwerpunkt auf den Verfahren der Nachrichtentechnik. Für Leser, die eine deutschsprachige Zweitdarstellung suchen.}

\lit{28}{Willems, W.: \textit{Codierungstheorie und Kryptographie}. Birkhäuser 2008.}{\textbf{Auf Deutsch und mathematisch streng.} Verbindet beide Hälften dieses Bandes in einem Band und entwickelt die algebraischen Grundlagen --- endliche Körper, Gruppen, elliptische Kurven --- von Grund auf. Die passende Vertiefung zu den Kapiteln~6 und 12 bis 14.}

\begin{merksatz}
\boxtitel{Wenn Sie nur drei Bücher kaufen}
\begin{center}
\renewcommand{\arraystretch}{1.6}
\begin{tabular}{@{}P{1.0cm}P{3.0cm}P{4.2cm}@{}}
\toprule
\textbf{[5]} & \textbf{Cover, Thomas} & \textbf{Für die Theorie}\\
\textbf{[13]} & \textbf{Katz, Lindell} & Für die Kryptographie\\
\textbf{[17]} & \textbf{MacKay} & \textbf{Für den Zusammenhang}\\
\bottomrule
\end{tabular}
\end{center}

\begin{center}
\begin{minipage}{0.88\textwidth}
\kursiv{Cover und Thomas enthalten die Teile I bis III mit allen Beweisen. Katz und Lindell behandeln Teil~IV und V mit der begrifflichen Strenge, die dieses Fach erst seit wenigen Jahrzehnten hat. MacKay steht nicht zufällig an dritter Stelle: Er ist der einzige, der Informationstheorie, Statistik und Lernverfahren als \kursiv{eine} Sache behandelt --- und er ist frei verfügbar.}
\end{minipage}
\end{center}

\begin{center}
\begin{minipage}{0.88\textwidth}
\kursiv{Vier frei verfügbare Ergänzungen: [24] Shannons Arbeit von 1948, die man gelesen haben sollte; [1] Anderson für alles, was über die Mathematik hinausgeht; [18] als Nachschlagewerk und [4] für Sicherheitsbeweise. Auf Deutsch: [3] Beutelspacher als Einstieg, [20] Paar und Pelzl für die Praxis, [28] Willems für die Algebra.}
\end{minipage}
\end{center}
\end{merksatz}

\bigskip
\begin{flushright}
\itshape\color{grau}
Information ist die\\
Auflösung von Ungewissheit.\\
Geheimhaltung ist\\
ihre Verweigerung.\\[10pt]
\normalfont\color{grau}\upshape
--- die beiden Hälften\\
\itshape dieses Bandes.
\end{flushright}
